\chapter{Les questions à poser à une source}\label{ch:grille}

Avant de transcrire une page, on pose à la source une dizaine de questions. Chacune a
coûté, sur l'un des fonds de ce livre, des jours de travail à qui ne l'avait pas posée.
Les réponses se trouvent presque toutes en une heure, avec un navigateur et la commande
\cmd{curl}.

\section{L'accès}

\begin{description}
  \item[Le fonds est-il servi entier ?] Une source qui ne montre qu'une sélection, ou qui
        renvoie pour le reste au site d'un partenaire, donne un inventaire troué. Le projet
        Gallica s'est donné une seule règle de sélection : un volume n'entre au catalogue
        que si Gallica le sert en entier.
  \item[Sous quelle forme ?] Un PDF par dossier (Montpellier, les archives Bourbaki, la
        Clay), des images par page derrière un protocole (IIIF chez Gallica et à la
        bibliothèque de Hanovre), des fichiers bruts sur un serveur (l'IHÉS), ou rien en
        ligne du tout (la BnF pour les papiers de Lasserre). Le format décide des scripts de
        copie, de la résolution disponible et de la façon d'afficher le fac-similé.
  \item[Le serveur se laisse-t-il afficher ailleurs ?] Deux en-têtes HTTP le disent.
        \texttt{X-Frame-Options} (ou \texttt{Content-Security-Policy: frame-ancestors})
        interdit d'afficher le document dans une page d'un autre site.
        \texttt{Access-Control-Allow-Origin} autorise, ou non, une page d'un autre site à
        lire les fichiers dans le navigateur. Sans ces autorisations, il faut un relais,
        c'est-à-dire un petit serveur à soi qui va chercher le document et le renvoie.
  \item[Le certificat est-il valide ?] Un certificat TLS expiré n'empêche pas un humain de
        cliquer sur « continuer », mais un cadre intégré dans une page reste blanc, sans
        message.
  \item[Combien de requêtes tolère-t-il ?] On ne le sait qu'en se faisant bloquer, et
        mieux vaut ne pas l'apprendre ainsi. Une pause entre deux requêtes, un plafond par
        heure, un \emph{user agent} qui nomme le projet, l'arrêt au premier refus.
\end{description}

\begin{terminal}
curl -sI https://exemple.org/document.pdf | grep -iE "x-frame|access-control|content-type"
\end{terminal}

\section{Le catalogue}

\begin{description}
  \item[Les cotes sont-elles stables ?] On cite la cote et l'adresse durable (un manifeste
        IIIF, une URL d'API), jamais un lien de visualiseur qui porte des paramètres de
        session.
  \item[Comment les pages sont-elles numérotées ?] Il y a souvent trois numérotations : celle
        de la page du fichier, celle que la bibliothèque a portée (foliotation), celle de
        l'auteur. On choisit celle qui pilote l'affichage, on note les autres quand on les
        lit, et on n'en déduit aucune.
  \item[Le catalogue date-t-il ?] Les datations du détenteur se citent telles quelles, entre
        leurs crochets s'il en met ; celles qu'on propose soi-même se marquent comme telles.
\end{description}

\section{Les droits}

\begin{description}
  \item[L'œuvre est-elle dans le domaine public ?] En France, les droits patrimoniaux
        courent soixante-dix ans après la mort de l'auteur ; le droit moral, dont le droit
        de divulgation d'une œuvre inédite, ne s'éteint pas. Grothendieck est mort en 2014,
        Quillen en 2011 : leurs notes sont protégées. Germain, Fourier ou Leibniz ne le sont
        plus.
  \item[Que permet le détenteur des images ?] Une numérisation a ses propres conditions
        (celles de Gallica, une \emph{Public Domain Mark}, un « © » sans plus). On les lit,
        on les date, on les cite.
  \item[Qui a déjà transcrit ?] Une édition imprimée, un projet de recherche, Wikisource. On
        ne recopie ni ne prend pour point de départ la transcription d'autrui sans son
        accord, et l'on renvoie à l'édition établie quand elle existe.
  \item[Y a-t-il des tiers ?] Les lettres reçues ont leurs propres auteurs. Les personnes
        privées nommées dans des papiers récents ont droit à leur vie privée.
\end{description}

\section{La main et la langue}

La dernière question est celle qu'on croit la première : peut-on lire ? Écriture,
abréviations, langue, notation de l'époque, figures, pages retournées ou écrites au dos
d'autres documents. Une machine lit mieux une dactylographie qu'un brouillon, et un
brouillon de 1960 mieux qu'un brouillon de 1680. Un échantillon de vingt pages, transcrit
et relu contre l'image, en dit plus que toute estimation.

\begin{regle}
Avant tout script de copie, on remplit la grille pour la source, dans une issue, avec la
date de chaque constat. L'issue sert ensuite de référence quand la source change.
\end{regle}
